\documentclass{homework}
% \usepackage{graphicx} % Required for inserting images
\usepackage{amsmath}
\usepackage{gensymb}
\usepackage{subfig}
\usepackage{float}

\class{ECEN 5427}
\title{Homework 2}
\author{Sean Jennings}
\date{March 6, 2025}

\begin{document}
\maketitle

\section{Problem 1}

\newcommand{\pu}{\text{ p.u.}}
\begin{enumerate}
    \item To compute the total power passing through the lines, first compute the total reactance $X$ as:
    \begin{equation*}
        \frac{1}{X} = \frac{1}{X_A} + \frac{1}{X_B}
    \end{equation*}
    or alternatively
    \begin{equation*}
        X = \frac{X_A X_B}{X_A + X_B} = \frac{0.1^2}{2 \cdot 0.1}  \pu = 0.05 \pu
    \end{equation*}
    and the power through each line is the following function of the total power $P=200$ MW
    \begin{equation*}
        P_A = \frac{X}{X_A} P = \frac{X_B}{X_A + X_B} P = \boxed{100 \text{ MW}}
    \end{equation*}
    and similarly
    \begin{equation*}
        P_B = \frac{X}{X_B} P = \frac{X_A}{X_A + X_B} P = \boxed{100 \text{ MW}}
    \end{equation*}

    \item The total reactance between bus 2 and 3 via bus 1 is the sum $X_{12} + X_{13} = $ 0.3 p.u.
    Since this equals the reactance of line 2-3, the power is divided equally between the two path,
    with 100 MW flowing through each line. The direction of the power flows are: 100 MW from bus 2 to bus 1,
    100 MW from bus 1 to bus 3, and 100 MW from bus 2 to bus 3.

\end{enumerate}

\section{Problem 2}
\renewcommand{\MW}{\text{ MW}}
\begin{enumerate}
    \item Let $X_{123}$ denote the total reactance between buses 1 and 3, via bus 2. I.e. $X_{123} = X_{12} + X_{23}$.
    Since the power flow is inversely proportional to the reactances, we have
    \begin{equation*}
        \frac{X_{123}}{X_{13}} = \frac{P_{13}}{P_{12}}.
    \end{equation*}
    Solving for $X_{123}$ gives
    \begin{equation*}
        X_{123} = \frac{P_{13}}{P_{12}} X_{13} = \frac{140\MW}{160 \MW} \cdot 0.5 \pu = 0.4375 \pu
    \end{equation*}
    so $X_{12}$ is
    \begin{equation*}
        X_{12} = X_{123} - X_{23} = (0.4375 - 0.3) \pu = \boxed{0.1375 \pu}
    \end{equation*}

    \item Since all reactances are equal, the path via bus 1 has twice the reactance as the direct
    line between bus 2 and 3, and therefore half as much power will flow through bus 1. That is,
    2/3 of the total power (333.3 MW) is transmitted on line 2-3 and 1/3 of the total power (166.7 MW) on each
    of lines 1-2 and 1-3.  This exceeds the capacity of line 1-2.

    The total amount that can be delivered is 3 times the capacity of line 1-2, or 450 MW.
\end{enumerate}

\section{Problem 3}
\begin{enumerate}
    \item In the unconstrained case, the nodal prices are all the same and the system-wide marginal cost
    is equal to the marginal cost of the marginal generator. The total load is $400 + 80 + 40 \MW = 520 \MW$.
    Generator $D$ (with the lowest cost) runs at capacity and generator $C$ covers the remaining 120 MW.
    Generator $C$ is the marginal generator, so system-wide and each nodal marginal cost is 10 \$/MWh.

    \item Generators $C$ and $D$ are at the same node, so we just need to super impose each load. 40 MW
    is deliver locally to bus 3, so no power flows through the line in this case. To transfer 80 MW from
    bus 3 to bus 2, we have a reactance of 0.5 p.u. via bus 1, and a reactance of 0.3 p.u. over branch 2-3.
    So the power $P_{ij}$ flowing from bus $i$ to bus $j$ is:
    \begin{align*}
        P_{32} &= \frac{0.5}{0.5 + 0.3} \cdot 80 \MW = 50 \MW \\
        P_{31} &= P_{12} = \frac{0.3}{0.5 + 0.3} \cdot 80 \MW = 30 \MW
    \end{align*}

    With just the 400 MW load active, we have:
    \begin{align*}
        P_{31} &= \frac{0.5}{0.5 + 0.3} \cdot 400 \MW = 250 \MW \\
        P_{32} &= P_{21} = \frac{0.3}{0.5 + 0.3} \cdot 400 \MW = 150 \MW
    \end{align*}


    Note the sign convention: $P_{12}$ indicates the flow from bus 1 to 2 whereas $P_{21}$ indicates the flow from 2 to 1, so
    $P_{12} = - P_{21}$.

    Summing over these two superimposed loads gives total powers of 
    \begin{align*}
        P_{32} &= (50 + 150) \MW = 200 \MW \\
        P_{31} &= (250 + 30) \MW = 280 \MW \\
        P_{21} &= (150 - 30) \MW = 120 \MW
    \end{align*}

    So branch 1-3 exceeds the security constraint.

    \item Since branch 1-2 is operating well below capacity and generator $A$ has a lower
    cost than generator $B$, we will first search for a feasible solution assuming that
    generator $D$ continues to run at capacity, while generator $A$ covers some of
    generator $C$'s portion of the load at bus 1.
    
    Five eighths (5/8) of the generation at bus 3 flows through branch 1-3, whereas
    $0.2 / (0.3 + 0.3 + 0.2) = 1/4$ of generation at bus 2 flows through branch 1-3. So
    for each MW $\Delta P$ of load transferred from generator $C$ to $A$, 
    we have $(1/4 - 5/8) \MW = -3/8 \MW = -0.375 \MW$ additional power flow through branch 1-3.
    That is,
    \begin{equation*}
        \Delta P_{13} = -\frac{3}{8} \cdot \Delta P
    \end{equation*}
    where $\Delta P_{13}$ indicates the change in flow on branch 1-3.
    
    Since we need 30 MW less power flow (i.e. $\Delta P_{13} = -30 \MW$) to relief the transmission constraint, we should take
    $\Delta P$ to be
    \begin{equation*}
        \Delta P = \frac{8}{3} \cdot 30 \MW = 80 \MW
    \end{equation*}

    Then, the new generation levels are:
    \begin{align*}
        P_A &= 80  \MW \\
        P_B &= 0 \MW \\
        P_C &= (120 - 80) \MW = 40 \MW \\
        P_D &= 400 \MW
    \end{align*}

    Let's check that this new generation profile is feasible. Now, the total power injected at bus 2 is 0 MW
    (load equals generation at bus 2). The total power injected at 
    bus 3 is $P_3 = 400 + 40 - 40 \MW = 400 \MW$ which results in the branch flows:
    \begin{align*}
        P_{31} &= \frac{0.5}{0.8} \cdot 400 \MW = 250 \MW \\
        P_{32} &= P_{21} = \frac{0.3}{0.8} \cdot 400 \MW = 150  \MW \\
    \end{align*}

    Branch 1-3 is now exactly at capacity which is what we'd expect from the optimal dispatch (since it is the solution to
    a linear program).
\end{enumerate}

\section{Bonus Problem 1}

Additional load at bus 3 is supplied by generator $C$ so the nodal price at bus 3 is 10 \$/MWh.
Load at bus 2 is supplied by generator $A$ giving a nodal price of 12 \$ / MWh.

Due to the security constraint, an additional load of $\Delta P_1$ at bus 1 is accomodated by increasing
the power from generator $A$ by $\Delta P_A > 0$ and decreasing power from generator $C$ by $\Delta P_C < 0$.
The power balance constraint is:
\begin{equation}
    \Delta P_1 = \Delta P_A + \Delta P_B
    \label{eq:power-balance}
\end{equation}
The change in the power flow $\Delta P_{12}$ over branch 1-2 is a weighted sum of $\Delta P_A$ and $\Delta P_B$
(with weights determined by the line impedances, as before):
\begin{equation*}
    \Delta P_{12} = 1/4 \Delta P_A + 5/8 \Delta P_C
\end{equation*}
To remain as close as possible without violating the security constraint, we set $\Delta P_{12} = 0$
which gives
\begin{equation}
    \Delta P_C = - \frac{2}{5} \Delta P_A
    \label{eq:impedance-frac}
\end{equation}

The total change in cost $\Delta C$ of additional generation is:
\begin{equation}
    \Delta C = p_A \Delta P_A + p_C \Delta P_C
    \label{eq:total-cost}
\end{equation}
where $p_A = 12$ \$/MWh and $p_C = 10$ \$/MWh.

Finally, the marginal cost at bus 1 is $\Delta C / \Delta P_1$, normalized by the amount of additional load.
Combining equations (\ref{eq:power-balance}), (\ref{eq:impedance-frac}), and (\ref{eq:total-cost}) gives
the marginal price
\begin{align*}
    \frac{C}{\Delta P_1} &= \frac{p_A \Delta P_A + p_C \Delta P_C}{\Delta P_A + \Delta P_C} \\
        &= \frac{(p_A - 2/5 p_C) \Delta P_A } {(1-2/5) \Delta P_A} \\
        &= \frac{5}{3} p_A - \frac{2}{3} p_C \\
        &= 13.33 \text{ \$/MWh}
\end{align*}

That this cost is less than the marginal cost of generator $B$ validates our decision in the previous
problem to transfer production from generator $C$ to generator $A$.

If generator $B$ had a marginal cost
of 13 \$ / MWh , it would be cheaper to start up generator $B$, since the amount of power $\Delta P_{CB}$ transferred
from $C$ to $B$ would be less in magnitude than $\Delta P_{CA}$. Even though the marginal cost of $B$ is higher than $A$,
the weighted sum along the transmission constraint is actually lower!


\section{Bonus Problem 2}

\begin{enumerate}
    \item When disconnected, the generation at each node must equal demand. That is, $P_A = 300$ MW and
    $P_B = 900$ MW. The marginal costs are therefore
    \begin{align*}
        MC_A = 9 + 0.015 P_A = 13.50 \text{ \$/MW} \\
        MC_B = 11 + 0.02 P_B = 29.00 \text{ \$/MW}
    \end{align*}

    \item With no transmission constraints, the marginal prices will equalize, giving
    a relation between the power produced by each generator:
    \begin{equation*}
        9 + 0.015 P_A = 11 + 0.02 P_B.
    \end{equation*}
    Solving for $P_A$:
    \begin{equation}
        P_A = \frac{1}{0.015} (2 + 0.02 P_B)
        \label{eq:solve-for-A}
    \end{equation}

    The power balance condition is that total generation meet total demand, i.e.
    \begin{equation}
        P_A + P_B = 1200 \MW
        \label{eq:power-balance2}
    \end{equation}
    Substituting equation (\ref{eq:solve-for-A}) into (\ref{eq:power-balance2})
    gives:
    \begin{align*}
        1200 &= \frac{1}{0.015} (2 + 0.02 P_B) + P_B \\
            &= 133.33 + 2.33 P_B.
    \end{align*}
    so the solution to the system of equations is
    \begin{align*}
        P_A &= 742.86 \MW \\
        P_B &= 457.14 \MW,
    \end{align*}

    the equilibrium price is:
    \begin{equation*}
        MC_A = MC_B = 20.14 \$ / \MW,
    \end{equation*}
    and power flow from bus $A$ to bus $B$ is
    \begin{equation*}
        P_{AB} = (742.86 - 300) \MW = (900 - 457.14) \MW = \boxed{442.86 \MW}
    \end{equation*}

    \item If transmission capacity is 350 MW, then the system will be transmission constrained, and the transmission line
    will run at capacity, with 350 MW flowing from bus $A$ to bus $B$.  Generator $A$ will produce $P_A = 650 MW$, and 
    generator $B$, $P_B = 550 MW$. The clearing prices will be
    \begin{align*}
        MC_A = 18.75 \text{ \$/MW} \\
        MC_B = 22.00 \text{ \$/MW}
    \end{align*}

    \item The revenues, $R_A$ and $R_B$, are just the generation amount times the marginal cost:
    \begin{align*}
        R_A &= P_A \cdot MC_A = 12,187.50 \text{ \$} \\
        R_B &= P_B \cdot MC_B = 12,100.00 \text{ \$}.
    \end{align*}



    


    
    





    
    
\end{enumerate}





\end{document}
